\documentclass[12pt,a4paper]{article}

% ── Packages ──────────────────────────────────────────────────────────
\usepackage[margin=1in]{geometry}
\usepackage{amsmath,amssymb,amsthm,mathtools}
\usepackage{hyperref}
\usepackage{cleveref}
\usepackage{tikz-cd}
\usepackage{tikz}
\usepackage{enumitem}
\usepackage{xcolor}
\usepackage{fancyhdr}
\usepackage{everypage}
\usepackage{abstract}
\usepackage{booktabs}
\usepackage{float}

% ── GrokRxiv DOI sidebar ─────────────────────────────────────────────
\definecolor{grokgray}{RGB}{110,110,110}
\AddEverypageHook{%
  \ifnum\value{page}=1
    \begin{tikzpicture}[remember picture, overlay]
      \node[rotate=90, anchor=south, font=\Large\sffamily\bfseries\color{grokgray}, inner sep=0pt]
      at ([xshift=38pt, yshift=0.52\paperheight]current page.south west)
      {GrokRxiv:2026.03.repetitive-elements-self-modifying\quad [\,bio.PL\,]\quad 31 Mar 2026};
    \end{tikzpicture}
  \fi
}

% ── Theorem environments ─────────────────────────────────────────────
\newtheorem{theorem}{Theorem}[section]
\newtheorem{lemma}[theorem]{Lemma}
\newtheorem{proposition}[theorem]{Proposition}
\newtheorem{corollary}[theorem]{Corollary}
\newtheorem{definition}[theorem]{Definition}
\newtheorem{example}[theorem]{Example}
\newtheorem{remark}[theorem]{Remark}
\newtheorem{axiom}[theorem]{Axiom}

% ── Page style ────────────────────────────────────────────────────────
\pagestyle{fancy}
\fancyhf{}
\fancyhead[L]{\small Repetitive Elements as Self-Modifying Code}
\fancyhead[R]{\small DNA-Lang Paper 5 of 7}
\fancyfoot[C]{\thepage}
\renewcommand{\headrulewidth}{0.4pt}

% ── Macros ────────────────────────────────────────────────────────────
\newcommand{\Genome}{\mathcal{G}}
\newcommand{\TE}{\mathsf{TE}}
\newcommand{\Repeat}{\mathsf{Repeat}}
\newcommand{\LINE}{\mathsf{LINE}}
\newcommand{\SINE}{\mathsf{SINE}}
\newcommand{\ERV}{\mathsf{ERV}}
\newcommand{\Sat}{\mathsf{Sat}}
\newcommand{\Hom}{\mathrm{Hom}}
\newcommand{\id}{\mathrm{id}}
\newcommand{\op}{\mathrm{op}}
\newcommand{\Cat}{\mathbf{Cat}}
\newcommand{\Set}{\mathbf{Set}}
\newcommand{\Seq}{\mathbf{Seq}}

\title{\textbf{Repetitive Elements as Self-Modifying Code:\\
A Type-Theoretic Framework for Genome Plasticity}}

\author{Matthew Long \\
\textit{The YonedaAI Collaboration} \\
Chicago, IL \\
\texttt{mlong@yonedaai.org}}

\date{March 31, 2026}

\begin{document}
\maketitle

\begin{abstract}
More than 45\% of the human genome consists of repetitive and transposable elements,
long dismissed as ``junk DNA.'' We present a type-theoretic framework that reconceptualizes
these elements as a sophisticated self-modifying code system within the DNA-Lang programming
language paradigm. DNA transposons are formalized as linear \emph{cut-and-paste} operations,
retrotransposons (LINEs and SINEs) as \emph{copy-and-paste} macros with dependency
relationships, satellite DNA as memory alignment primitives, and endogenous retroviruses
as legacy library imports. We introduce the $\Repeat\langle\mathsf{SelfModifying}\rangle$
type, prove that transposition defines an endofunctor on the genome category, establish
a G\"odelian limitation on transposon self-prediction, and formalize transposon
\emph{domestication} as type naturalization. Our Haskell implementation provides executable
models of all constructs. This is Paper~5 of~7 in the DNA-Lang series, building on the
modular framework established in Papers~1--4 and providing foundations for developmental
programs (Paper~6) and structural DNA (Paper~7).
\end{abstract}

\tableofcontents
\newpage

% ══════════════════════════════════════════════════════════════════════
\section{Introduction}\label{sec:intro}
% ══════════════════════════════════════════════════════════════════════

The human genome comprises approximately 3.2 billion base pairs, yet fewer than 2\%
encode proteins. The remaining landscape is dominated by repetitive elements: DNA
transposons, retrotransposons (both autonomous LINEs and non-autonomous SINEs),
satellite DNA, and endogenous retroviruses (ERVs). Together, these constitute over
45\% of the genome~\cite{lander2001,debarros2008}.

For decades, these elements were labelled ``selfish DNA'' or ``junk DNA,'' understood
merely as genomic parasites that replicate at the host's expense~\cite{orgel1980,doolittle1980}.
This view has undergone a profound revision. Transposable elements (TEs) serve as
reservoirs of regulatory innovation~\cite{feschotte2008}, drive genome
restructuring~\cite{kazazian2004}, and have been repeatedly \emph{domesticated} for
host function---most dramatically in the case of V(D)J recombination, the engine of
adaptive immunity~\cite{agrawal1998}.

In the DNA-Lang framework, we interpret the genome as a programming language whose source
code is simultaneously its own runtime. Repetitive elements, in this view, are not inert
detritus but a \emph{self-modifying code system}---instructions that rewrite the very
program they inhabit. This places genomes alongside the classical self-modifying
architectures studied in computer science, from von~Neumann's stored-program concept to
Lisp macros and JIT compilers.

\begin{definition}[Self-Modifying Genome]\label{def:selfmod}
A \emph{self-modifying genome} is a triple $(\Genome, \TE, \tau)$ where:
\begin{enumerate}[label=(\roman*)]
  \item $\Genome$ is a finite sequence over the alphabet $\{A, T, G, C\}$,
  \item $\TE \subseteq \Genome$ is the set of transposable element loci,
  \item $\tau : \TE \times \Genome \to \Genome$ is the transposition function satisfying
    $|\tau(t,g)| \geq |g|$ for retrotransposons (copy semantics) and
    $|\tau(t,g)| = |g|$ for DNA transposons (move semantics).
\end{enumerate}
\end{definition}

The contributions of this paper are:
\begin{enumerate}[label=\arabic*.]
  \item A formal type system for repetitive elements, the $\Repeat\langle\mathsf{SelfModifying}\rangle$ type, capturing transposition mode, activity state, and host impact (\cref{sec:repeat-type}).
  \item Proof that transposition defines an endofunctor on the genome category (\cref{sec:endofunctor}).
  \item Formalization of LINE/SINE dependency as a dependent type relationship (\cref{sec:sines}).
  \item A G\"odelian impossibility result on transposon self-prediction (\cref{sec:goedel}).
  \item A complete Haskell implementation (\cref{sec:implementation}).
\end{enumerate}

\subsection{Relationship to the DNA-Lang Series}

This paper is the fifth in a modular series of seven:
\begin{enumerate}[label=\textbf{Paper \arabic*:}]
  \item \textit{Regulatory Sequences as Higher-Order Functions}
  \item \textit{Coding Sequences as Executable Functions}
  \item \textit{Non-Coding RNA as Type-Level Programming}
  \item \textit{Epigenetic Marks as Compiler Pragmas}
  \item \textit{Repetitive Elements as Self-Modifying Code} (this paper)
  \item \textit{Developmental Programs as Build Systems}
  \item \textit{Structural DNA as Memory Architecture}
\end{enumerate}

Each paper in the modular framework builds upon previous ones through hierarchical
composition. Repetitive elements interact with regulatory sequences (Paper~1) by
providing novel promoters upon insertion, with coding sequences (Paper~2) through
exon shuffling and gene duplication, with non-coding RNA (Paper~3) as SINEs generate
functional RNA transcripts, and with epigenetic marks (Paper~4) as chromatin state
controls TE activity.


% ══════════════════════════════════════════════════════════════════════
\section{DNA Transposons as Cut-and-Paste}\label{sec:cut-paste}
% ══════════════════════════════════════════════════════════════════════

DNA transposons, also known as Class~II transposable elements, mobilize via a
\emph{cut-and-paste} mechanism. The transposase enzyme recognizes terminal inverted
repeats (TIRs) flanking the element, excises it from the donor site, and reinserts
it at a target site. This is analogous to \emph{move semantics} in programming languages:
the element exists at exactly one location at any time.

\begin{definition}[Cut-and-Paste Transposition]\label{def:cut-paste}
Let $\Genome = \alpha \cdot t \cdot \beta$ where $t \in \TE$ is a DNA transposon
flanked by sequences $\alpha$ and $\beta$. A cut-and-paste transposition to target
site $i$ in $\beta = \beta_1 \cdot \beta_2$ (with $|\beta_1| = i$) is:
\[
  \tau_{\mathrm{cut}}(t, \Genome) = \alpha \cdot \beta_1 \cdot t \cdot \beta_2
\]
satisfying the \emph{conservation law}: $|\tau_{\mathrm{cut}}(t, \Genome)| = |\Genome|$.
\end{definition}

This conservation law is the genomic analogue of linear type systems in programming
languages, where resources must be used exactly once.

\begin{theorem}[Linear Type Preservation]\label{thm:linear}
Cut-and-paste transposition preserves the linear type invariant. If $\Genome$ contains
exactly $n$ copies of transposon $t$, then $\tau_{\mathrm{cut}}(t, \Genome)$ also
contains exactly $n$ copies.
\end{theorem}

\begin{proof}
By \cref{def:cut-paste}, $\tau_{\mathrm{cut}}$ excises $t$ from position $|\alpha|$
and reinserts it at position $|\alpha| + |\beta_1|$. The multiset of subsequences is
preserved since the operation is a permutation of the genome's factorization. Hence the
count of any subsequence, including $t$, is invariant.
\end{proof}

\subsection{Terminal Inverted Repeats as Delimiters}

TIRs serve as the syntactic delimiters that the transposase recognizes. They form a
palindromic boundary:

\begin{definition}[Terminal Inverted Repeat]\label{def:tir}
A \emph{terminal inverted repeat} for element $t$ is a pair $(l, r)$ of sequences
satisfying $r = \overline{l^R}$, where $\overline{(\cdot)}$ denotes complement and
$(\cdot)^R$ denotes reversal. The element is bounded as $l \cdot t_{\mathrm{body}} \cdot r$.
\end{definition}

This is precisely a \emph{matched delimiter} in the syntactic sense: the transposase
parser recognizes the opening TIR, scans for the matching closing TIR, and the enclosed
body is the transposon payload.

\subsection{Target Site Duplication}

Upon insertion, the transposase creates a staggered cut at the target site, producing
short direct repeats flanking the inserted element---the \emph{target site duplication}
(TSD). This is the ``footprint'' left by transposition:

\begin{proposition}[TSD as Insertion Trace]\label{prop:tsd}
Every cut-and-paste insertion at a target site $s$ of length $k$ produces a genome
$\alpha \cdot s \cdot \mathrm{TIR}_L \cdot t_{\mathrm{body}} \cdot \mathrm{TIR}_R
\cdot s \cdot \beta$, where $s$ is duplicated. The corrected conservation law is:
\[
  |\tau_{\mathrm{cut}}(t, \Genome)| = |\Genome| + |s|
\]
accounting for the TSD.
\end{proposition}

\subsection{The Transposase as an Interpreter}

The transposase protein is the \emph{interpreter} of the transposon instruction.
It reads the TIR delimiters, parses the transposon body, and executes the
cut-and-paste operation. In the DNA-Lang type system:

\[
  \mathsf{transposase} : \mathrm{TIR} \to \Genome \to (\Genome, \mathsf{InsertionSite})
\]

The transposase is itself encoded within many DNA transposons, making these elements
\emph{self-interpreting code}: programs that contain their own interpreter.


% ══════════════════════════════════════════════════════════════════════
\section{Retrotransposons as Copy-and-Paste}\label{sec:copy-paste}
% ══════════════════════════════════════════════════════════════════════

Class~I elements, or retrotransposons, mobilize through an RNA intermediate: the
element is transcribed to RNA, reverse-transcribed back to DNA, and the new copy
is inserted at a target site. The original copy remains, yielding \emph{copy semantics}.

\begin{definition}[Copy-and-Paste Transposition]\label{def:copy-paste}
Let $\Genome = \alpha \cdot t \cdot \beta$ where $t$ is a retrotransposon. A
copy-and-paste transposition to target site $i$ in $\beta = \beta_1 \cdot \beta_2$ is:
\[
  \tau_{\mathrm{copy}}(t, \Genome) = \alpha \cdot t \cdot \beta_1 \cdot t' \cdot \beta_2
\]
where $t'$ is a (possibly mutated) copy of $t$. The \emph{expansion law}:
$|\tau_{\mathrm{copy}}(t, \Genome)| = |\Genome| + |t'|$.
\end{definition}

\begin{theorem}[Monotonic Genome Growth]\label{thm:growth}
Under repeated copy-and-paste transposition, the genome size is monotonically
non-decreasing:
\[
  |\Genome_0| \leq |\Genome_1| \leq |\Genome_2| \leq \cdots
\]
Moreover, without deletion mechanisms, $|\Genome_n| = |\Genome_0| + \sum_{i=0}^{n-1} |t_i'|$.
\end{theorem}

\begin{proof}
Immediate from \cref{def:copy-paste}: each copy-and-paste event adds $|t'| > 0$ bases.
\end{proof}

This explains the empirical observation that retrotransposon-rich genomes tend to be
large: the maize genome, for instance, has roughly doubled in size over the past
3 million years due to retrotransposon expansion~\cite{sanmiguel1998}.

\subsection{LINEs: Long Interspersed Nuclear Elements}

LINEs are autonomous retrotransposons encoding their own reverse transcriptase (RT)
and endonuclease. The human L1 element, at approximately 6~kb, is the dominant LINE,
with over 500{,}000 copies comprising roughly 17\% of the human genome.

\begin{definition}[LINE Structure]\label{def:line}
A LINE $\ell$ has the structure:
\[
  \ell = \mathsf{5'UTR} \cdot \mathsf{ORF1} \cdot \mathsf{ORF2} \cdot \mathsf{3'UTR} \cdot \mathsf{polyA}
\]
where $\mathsf{ORF1}$ encodes an RNA-binding protein and $\mathsf{ORF2}$ encodes
a bifunctional protein with endonuclease and reverse transcriptase domains.
\end{definition}

The L1 retrotransposition cycle implements \emph{target-primed reverse transcription}
(TPRT):
\begin{enumerate}[label=\arabic*.]
  \item Transcription: $\ell_{\mathrm{DNA}} \xrightarrow{\mathrm{Pol\,II}} \ell_{\mathrm{RNA}}$
  \item Translation: $\ell_{\mathrm{RNA}} \xrightarrow{\mathrm{ribosome}} (\mathsf{ORF1p}, \mathsf{ORF2p})$
  \item Target nicking: $\mathsf{ORF2p}_{\mathrm{EN}} : \Genome \to (\Genome, \mathsf{Nick})$
  \item Reverse transcription: $\mathsf{ORF2p}_{\mathrm{RT}} : \ell_{\mathrm{RNA}} \to \ell_{\mathrm{cDNA}}$
  \item Integration: $\ell_{\mathrm{cDNA}} \xrightarrow{\mathrm{repair}} \Genome'$
\end{enumerate}

\subsection{SINEs: Short Interspersed Nuclear Elements}

SINEs are non-autonomous: they do not encode their own transposition machinery but
\emph{parasitize} the LINE machinery. The most abundant SINE in humans is the Alu
element (approximately 300~bp), with over 1.1~million copies comprising roughly 11\%
of the genome.

\begin{definition}[SINE Dependency]\label{def:sine-dep}
A SINE $s$ is a retrotransposon satisfying:
\begin{enumerate}[label=(\roman*)]
  \item $s$ does not encode reverse transcriptase or endonuclease,
  \item $s$ contains a Pol~III internal promoter,
  \item Transposition of $s$ requires $\exists \ell \in \LINE.\; \mathsf{active}(\ell)$.
\end{enumerate}
\end{definition}

This dependency relationship will be formalized categorically in \cref{sec:sines}.


% ══════════════════════════════════════════════════════════════════════
\section{The $\Repeat\langle\mathsf{SelfModifying}\rangle$ Type}\label{sec:repeat-type}
% ══════════════════════════════════════════════════════════════════════

We now introduce the central type of our framework.

\begin{definition}[$\Repeat\langle\mathsf{SelfModifying}\rangle$ Type]\label{def:repeat-type}
The type $\Repeat\langle S \rangle$ is a parameterized algebraic data type:
\begin{align*}
\Repeat\langle S \rangle \;::=\;&\; \mathsf{DNATransposon}\{\\
  &\quad \mathsf{tirs} : \mathrm{TIR} \times \mathrm{TIR},\\
  &\quad \mathsf{transposase} : \mathrm{Gene},\\
  &\quad \mathsf{mode} : \mathsf{CutPaste},\\
  &\quad \mathsf{activity} : S \}\\
  \mid\;&\; \mathsf{Retrotransposon}\{\\
  &\quad \mathsf{class} : \LINE \mid \SINE,\\
  &\quad \mathsf{rt} : \mathrm{Maybe}(\mathrm{ReverseTranscriptase}),\\
  &\quad \mathsf{mode} : \mathsf{CopyPaste},\\
  &\quad \mathsf{activity} : S \}\\
  \mid\;&\; \mathsf{SatelliteDNA}\{\\
  &\quad \mathsf{unit} : \mathrm{Sequence},\\
  &\quad \mathsf{copies} : \mathbb{N},\\
  &\quad \mathsf{location} : \mathsf{Centromere} \mid \mathsf{Telomere} \mid \mathsf{Pericentromeric} \}\\
  \mid\;&\; \ERV\{\\
  &\quad \mathsf{ltrs} : \mathrm{LTR} \times \mathrm{LTR},\\
  &\quad \mathsf{genes} : [\mathrm{Gene}],\\
  &\quad \mathsf{intact} : \mathrm{Bool},\\
  &\quad \mathsf{activity} : S \}
\end{align*}
where the type parameter $S$ tracks the activity state, drawn from the lattice:
\[
  \mathsf{Active} \sqsubset \mathsf{Suppressed} \sqsubset \mathsf{Fossilized} \sqsubset \mathsf{Domesticated}
\]
\end{definition}

\begin{remark}
The activity lattice reflects the evolutionary trajectory of transposable elements:
active elements transpose freely; suppressed elements are silenced by host mechanisms
(DNA methylation, piRNA); fossilized elements have accumulated inactivating mutations;
domesticated elements have been co-opted for host function.
\end{remark}

\begin{definition}[Host Impact Functor]\label{def:impact}
The \emph{host impact} of a repeat element is captured by the functor
$\mathsf{Impact} : \Repeat \to \mathsf{Effect}$ where:
\[
\mathsf{Effect} ::= \mathsf{Neutral} \mid \mathsf{Deleterious}(\mathbb{R}^+)
  \mid \mathsf{Regulatory} \mid \mathsf{Structural} \mid \mathsf{Exapted}(\mathrm{Function})
\]
\end{definition}


% ══════════════════════════════════════════════════════════════════════
\section{Transposition as Endofunctor}\label{sec:endofunctor}
% ══════════════════════════════════════════════════════════════════════

We now establish the central categorical result: transposition defines an endofunctor
on the genome category.

\begin{definition}[Genome Category]\label{def:genome-cat}
The \emph{genome category} $\mathcal{C}_\Genome$ has:
\begin{itemize}
  \item \textbf{Objects:} Genome states $\Genome_i$, each a finite sequence over $\{A,T,G,C\}$
    with annotated loci.
  \item \textbf{Morphisms:} Genomic transformations $f : \Genome_i \to \Genome_j$, including
    point mutations, insertions, deletions, and rearrangements.
  \item \textbf{Composition:} Sequential application of transformations.
  \item \textbf{Identity:} The null transformation $\id_\Genome$.
\end{itemize}
\end{definition}

\begin{theorem}[Transposition Endofunctor]\label{thm:endofunctor}
The transposition operation $T : \mathcal{C}_\Genome \to \mathcal{C}_\Genome$ is an
endofunctor, where:
\begin{enumerate}[label=(\roman*)]
  \item On objects: $T(\Genome_i) = \tau(t, \Genome_i)$ for a fixed active element $t$.
  \item On morphisms: $T(f : \Genome_i \to \Genome_j) = f' : T(\Genome_i) \to T(\Genome_j)$
    where $f'$ applies $f$ to all positions not occupied by the transposed element, and
    adjusts coordinates accordingly.
\end{enumerate}
\end{theorem}

\begin{proof}
We verify the functor laws.

\textbf{Identity preservation:} $T(\id_{\Genome_i})$ must equal $\id_{T(\Genome_i)}$.
The identity transformation changes no positions; after transposition of $t$, the identity
on the resulting genome still changes no positions. Hence $T(\id) = \id_{T(\Genome)}$.

\textbf{Composition preservation:} We must show $T(g \circ f) = T(g) \circ T(f)$.
Let $f : \Genome_i \to \Genome_j$ and $g : \Genome_j \to \Genome_k$. The transposition
of $t$ commutes with mutations at non-overlapping positions. For the case where $f$ or
$g$ affects the transposon locus, the coordinate adjustment in the definition of $T$
on morphisms ensures that
\[
  T(g \circ f)(\Genome_i) = T(g)(T(f)(\Genome_i))
\]
by the associativity of sequence operations and the determinism of coordinate remapping.
\end{proof}

\begin{corollary}[Repeated Transposition as Iteration]\label{cor:iteration}
For an active transposon $t$, the $n$-fold transposition $T^n = \underbrace{T \circ \cdots \circ T}_{n}$
is also an endofunctor, and the collection $\{T^n \mid n \in \mathbb{N}\}$ forms a
monoid in the endofunctor category $\mathrm{End}(\mathcal{C}_\Genome)$.
\end{corollary}

The following diagram commutes for any genomic mutation $f$:
\[
\begin{tikzcd}
  \Genome_i \arrow[r, "f"] \arrow[d, "T"'] & \Genome_j \arrow[d, "T"] \\
  T(\Genome_i) \arrow[r, "T(f)"'] & T(\Genome_j)
\end{tikzcd}
\]

\begin{proposition}[Natural Transformation of Transposons]\label{prop:nat-trans}
Given two active transposons $t_1, t_2 \in \TE$, the interaction between their
respective endofunctors $T_1, T_2$ is a natural transformation
$\eta : T_1 \Rightarrow T_2$ if and only if $t_1$ and $t_2$ share the same target
site preference distribution.
\end{proposition}


% ══════════════════════════════════════════════════════════════════════
\section{LINEs as Autonomous Macros}\label{sec:lines}
% ══════════════════════════════════════════════════════════════════════

We formalize LINEs as \emph{autonomous macros}---self-contained expansion units
that carry their own expansion machinery.

\begin{definition}[Autonomous Macro]\label{def:auto-macro}
An \emph{autonomous macro} is a pair $(t, \mathsf{expand}_t)$ where $t$ is a
transposable element and $\mathsf{expand}_t : t \to \Genome \to \Genome$ is the
expansion function encoded within $t$ itself. For LINEs:
\[
  \mathsf{expand}_\LINE = \lambda \ell.\; \lambda g.\;
    \mathsf{insert}(\mathsf{TPRT}(\mathsf{transcribe}(\ell)), g)
\]
\end{definition}

The L1 element is the paradigmatic autonomous macro:

\begin{proposition}[L1 Self-Sufficiency]\label{prop:l1}
The L1 LINE element is \emph{self-sufficient} in the sense that:
\[
  \forall g \in \Genome.\; \mathsf{expand}_{\mathrm{L1}}(\mathrm{L1}, g) \neq \bot
\]
provided the cellular transcription and translation machinery is available. That is,
L1 requires only the host's basic gene expression infrastructure, not any
TE-specific factors.
\end{proposition}

\subsection{L1 as a Monad}

The L1 retrotransposition cycle has monadic structure:

\begin{theorem}[L1 Monad]\label{thm:l1-monad}
Define the type constructor $L : \Genome \to \Genome$ by $L(g) = g$ augmented with a
fresh L1 insertion. Then $(L, \eta, \mu)$ is a monad where:
\begin{itemize}
  \item $\eta_g : g \to L(g)$ is the insertion of a single L1 copy (unit),
  \item $\mu_g : L(L(g)) \to L(g)$ flattens nested transposition events (join).
\end{itemize}
\end{theorem}

\begin{proof}
We verify the monad laws.

\textbf{Left unit:} $\mu \circ L(\eta) = \id_L$.
Applying $\eta$ inside $L$ creates a nested insertion; $\mu$ flattens it to a single
insertion, recovering the original $L(g)$.

\textbf{Right unit:} $\mu \circ \eta_L = \id_L$.
Wrapping $L(g)$ in a new $L$ layer and flattening returns $L(g)$.

\textbf{Associativity:} $\mu \circ L(\mu) = \mu \circ \mu_L$.
Both sides flatten three levels of nesting to one. Since insertion positions are
determined by the TPRT mechanism (which is position-independent at the categorical
level), the order of flattening does not matter.
\end{proof}

\subsection{5' Truncation as Lossy Compression}

Most L1 copies in the genome are 5'-truncated: the TPRT mechanism begins at the 3' end,
and premature termination leaves incomplete copies. Of the approximately 500{,}000
L1 copies in the human genome, fewer than 100 are full-length and retrotransposition-competent.

\begin{definition}[5' Truncation]\label{def:truncation}
The \emph{truncation function} $\mathsf{trunc}_k : \LINE \to \LINE$ removes the first
$k$ bases:
\[
  \mathsf{trunc}_k(\mathsf{5'UTR} \cdot \mathsf{ORF1} \cdot \mathsf{ORF2} \cdot \mathsf{3'UTR} \cdot \mathsf{polyA})
  = \mathsf{suffix}_k
\]
A truncated LINE with $k > |\mathsf{5'UTR}| + |\mathsf{ORF1}|$ loses ORF1, rendering
it non-autonomous.
\end{definition}


% ══════════════════════════════════════════════════════════════════════
\section{SINEs as Dependent Macros}\label{sec:sines}
% ══════════════════════════════════════════════════════════════════════

SINEs represent a fundamentally different computational pattern: they are
\emph{dependent macros} that require external machinery for expansion.

\begin{definition}[Dependent Macro]\label{def:dep-macro}
A \emph{dependent macro} is a pair $(s, \pi)$ where $s$ is a transposable element
and $\pi : \LINE \to \mathsf{ExpansionMachinery}$ is a projection that extracts the
required machinery from an autonomous element. The expansion of $s$ is:
\[
  \mathsf{expand}_\SINE : \SINE \to \Pi(\ell : \LINE).\; \mathsf{active}(\ell) \to \Genome \to \Genome
\]
This is a \emph{dependent function type}: the SINE's expansion depends on the existence
and activity of a LINE element.
\end{definition}

\begin{theorem}[SINE--LINE Dependency as Dependent Type]\label{thm:sine-dep}
The relationship between SINEs and LINEs is captured by the dependent type:
\[
  \mathsf{SINETransposition} : \Sigma(\ell : \LINE).\; \mathsf{active}(\ell) \times (\SINE \to \Genome \to \Genome)
\]
That is, a SINE transposition event is a \emph{witness} to the existence of an active
LINE together with the transposition function.
\end{theorem}

\begin{proof}
By \cref{def:sine-dep}, every SINE transposition requires active LINE machinery.
The $\Sigma$-type packages this existential witness: to construct a
$\mathsf{SINETransposition}$, one must provide a specific LINE $\ell$, a proof of
$\mathsf{active}(\ell)$, and the resulting transposition function. This is precisely
the elimination rule for dependent pairs.
\end{proof}

\subsection{Alu Elements: The Most Successful Dependent Macro}

Alu elements are dimeric SINEs derived from the 7SL RNA gene. With over 1.1~million
copies, they are the most successful mobile element in the human genome by copy number.

\begin{example}[Alu Dependency Chain]\label{ex:alu}
The Alu retrotransposition dependency chain is:
\[
\begin{tikzcd}[column sep=large]
  \mathsf{Alu}_{\mathrm{RNA}} \arrow[r, "\text{hijacks}"] &
  \mathsf{L1\;ORF2p} \arrow[r, "\text{TPRT}"] &
  \mathsf{Alu}_{\mathrm{cDNA}} \arrow[r, "\text{insert}"] &
  \Genome'
\end{tikzcd}
\]
Alu cannot perform any of these steps autonomously; it is a pure \emph{client} of
L1 services.
\end{example}

\begin{proposition}[Alu Expansion Rate Bounded by L1 Activity]\label{prop:alu-bound}
The rate of Alu transposition $r_{\mathrm{Alu}}$ is bounded above by the L1 activity:
\[
  r_{\mathrm{Alu}} \leq k \cdot \sum_{\ell \in \LINE_{\mathrm{active}}} \mathsf{activity}(\ell)
\]
for some constant $k > 0$ depending on Alu copy number and transcriptional efficiency.
\end{proposition}


% ══════════════════════════════════════════════════════════════════════
\section{Satellite DNA as Memory Alignment}\label{sec:satellite}
% ══════════════════════════════════════════════════════════════════════

Satellite DNA consists of tandemly repeated short sequences (2--200~bp units) found
primarily at centromeres and telomeres. In the DNA-Lang framework, these serve as
\emph{memory alignment and padding primitives}.

\begin{definition}[Tandem Repeat]\label{def:tandem}
A \emph{tandem repeat} is a sequence $r^n = \underbrace{r \cdot r \cdots r}_{n}$
where $r$ is the repeat unit and $n$ is the copy number. The total length is $n \cdot |r|$.
\end{definition}

\begin{definition}[Satellite Classes]\label{def:sat-classes}
Satellite DNA is classified by repeat unit length:
\begin{center}
\begin{tabular}{@{}llll@{}}
\toprule
\textbf{Class} & \textbf{Unit Length} & \textbf{Location} & \textbf{CS Analogue} \\
\midrule
Satellite & 100--200 bp & Pericentromeric & Page alignment \\
Minisatellite & 10--100 bp & Subtelomeric & Cache-line alignment \\
Microsatellite & 2--9 bp & Dispersed & Word alignment \\
\bottomrule
\end{tabular}
\end{center}
\end{definition}

\subsection{Centromeric Satellites as Struct Padding}

The centromere, essential for chromosome segregation during cell division, is built upon
vast arrays of alpha-satellite DNA (171~bp repeat units). This is analogous to
\emph{struct padding} in C: the functional requirement (kinetochore assembly) demands a
specific memory alignment that is achieved by repeating a unit until the required size
is met.

\begin{theorem}[Centromeric Alignment]\label{thm:centromere}
The centromeric satellite array of length $L = n \cdot 171$ bp ensures that:
\begin{enumerate}[label=(\roman*)]
  \item The kinetochore binding region is aligned to a higher-order repeat (HOR)
    boundary of $m \cdot 171$ bp, where $m$ is the HOR period.
  \item The array length $L$ exceeds a minimum threshold $L_{\min}$ required for
    stable kinetochore assembly.
  \item The CENP-B box motif occurs at regular intervals within the array, providing
    ``alignment markers.''
\end{enumerate}
\end{theorem}

\subsection{Telomeric Repeats as Sentinel Values}

Telomeres consist of TTAGGG repeats (in vertebrates) and serve as \emph{sentinel values}
marking the ends of chromosomes. They prevent end-to-end fusion and degradation, analogous
to null terminators in C strings or sentinel nodes in linked lists.

\begin{proposition}[Telomere as Sentinel]\label{prop:telomere}
The telomeric repeat array $(\mathrm{TTAGGG})^n$ functions as a sentinel by:
\begin{enumerate}[label=(\roman*)]
  \item Preventing exonuclease degradation of coding sequence (buffer overflow protection),
  \item Signaling the physical end of the chromosome to the replication machinery,
  \item Providing a substrate for telomerase-mediated length maintenance (garbage collection).
\end{enumerate}
\end{proposition}


% ══════════════════════════════════════════════════════════════════════
\section{Endogenous Retroviruses as Legacy Imports}\label{sec:ervs}
% ══════════════════════════════════════════════════════════════════════

Endogenous retroviruses (ERVs) are remnants of ancient retroviral infections that
became fixed in the germline. They constitute approximately 8\% of the human genome.
In the DNA-Lang framework, these are \emph{legacy library imports}: foreign code
packages that were integrated into the codebase long ago and remain structurally
embedded even though most are no longer functional.

\begin{definition}[Endogenous Retrovirus]\label{def:erv}
An ERV $e$ has the structure:
\[
  e = \mathrm{LTR}_5 \cdot \mathsf{gag} \cdot \mathsf{pol} \cdot \mathsf{env} \cdot \mathrm{LTR}_3
\]
where $\mathsf{gag}$ encodes capsid proteins, $\mathsf{pol}$ encodes reverse
transcriptase and integrase, and $\mathsf{env}$ encodes envelope proteins. The
flanking long terminal repeats (LTRs) contain promoter and enhancer sequences.
\end{definition}

\begin{proposition}[ERV Decay]\label{prop:erv-decay}
Over evolutionary time, ERVs undergo systematic degradation:
\begin{enumerate}[label=(\roman*)]
  \item Accumulation of inactivating mutations in $\mathsf{gag}$, $\mathsf{pol}$, $\mathsf{env}$,
  \item Deletion of internal sequence by recombination between flanking LTRs, leaving
    solo LTRs,
  \item Epigenetic silencing via DNA methylation and heterochromatin.
\end{enumerate}
The human genome contains approximately 203{,}000 solo LTRs versus 9{,}000 intact ERV
loci---a ratio reflecting extensive recombinational deletion.
\end{proposition}

\subsection{Solo LTRs as Orphaned Headers}

When the internal genes of an ERV are deleted by homologous recombination between the
two LTRs, a single ``solo LTR'' remains. This is the genomic equivalent of an
\emph{orphaned header file}: the interface declaration persists (the LTR with its
promoter/enhancer signals) even though the implementation (the viral genes) has been
deleted.

\begin{theorem}[Solo LTR Regulatory Potential]\label{thm:solo-ltr}
A solo LTR retains transcription factor binding sites from the original viral LTR and
can function as:
\begin{enumerate}[label=(\roman*)]
  \item An alternative promoter for nearby host genes,
  \item An enhancer element acting at a distance,
  \item A source of novel regulatory sequence upon mutation.
\end{enumerate}
The regulatory potential is a function $\mathsf{RegPot} : \mathrm{LTR} \to \mathcal{P}(\mathrm{TF})$
mapping each solo LTR to the set of transcription factors it can bind.
\end{theorem}


% ══════════════════════════════════════════════════════════════════════
\section{Transposon Domestication as Naturalization}\label{sec:domestication}
% ══════════════════════════════════════════════════════════════════════

The most remarkable fate of a transposable element is \emph{domestication}: the
evolutionary process by which a parasitic element is co-opted for host function.
In the DNA-Lang framework, this is \emph{type naturalization}---the conversion of
foreign code into a native component.

\begin{definition}[Type Naturalization]\label{def:naturalization}
A \emph{naturalization} of a repeat element $r \in \Repeat\langle\mathsf{Active}\rangle$
is a type-changing transformation:
\[
  \mathsf{naturalize} : \Repeat\langle\mathsf{Active}\rangle \to \mathsf{HostGene}
\]
satisfying:
\begin{enumerate}[label=(\roman*)]
  \item \textbf{Loss of mobility:} $\mathsf{canTranspose}(\mathsf{naturalize}(r)) = \mathsf{False}$,
  \item \textbf{Gain of host function:} $\exists f.\; \mathsf{hostFunction}(\mathsf{naturalize}(r)) = \mathsf{Just}(f)$,
  \item \textbf{Selective constraint:} The naturalized element is under purifying selection
    ($dN/dS < 1$).
\end{enumerate}
\end{definition}

\subsection{V(D)J Recombination: The RAG Transposase}

The most celebrated case of transposon domestication is the RAG1/RAG2 recombinase
system that drives V(D)J recombination in the adaptive immune system of jawed
vertebrates~\cite{agrawal1998,hiom1998}.

\begin{theorem}[RAG as Domesticated Transposase]\label{thm:rag}
The RAG1 protein is a domesticated transposase satisfying:
\begin{enumerate}[label=(\roman*)]
  \item RAG1 contains a DDE catalytic triad characteristic of transposases,
  \item RAG1/RAG2 catalyze a cut-and-paste reaction on recombination signal
    sequences (RSSs) that are structurally analogous to TIRs,
  \item The reaction mechanism (synapsis, cleavage, hairpin formation) recapitulates
    DNA transposition,
  \item RAG1/RAG2 can catalyze \emph{in vitro} transposition, demonstrating retained
    ancestral activity.
\end{enumerate}
Therefore:
\[
  \mathsf{RAG} = \mathsf{naturalize}(\mathsf{Transib}_{\mathrm{transposase}})
\]
where $\mathsf{Transib}$ is the ancestral DNA transposon superfamily.
\end{theorem}

\subsection{Syncytins: Domesticated Retroviral Envelope Proteins}

The syncytin proteins, essential for placental development in mammals, are domesticated
ERV envelope ($\mathsf{env}$) genes~\cite{mi2000}.

\begin{example}[Syncytin Naturalization]\label{ex:syncytin}
\[
\begin{tikzcd}
  \ERV_{\mathsf{env}} \arrow[r, "\text{mutation}"] \arrow[d, "\text{purifying selection}"'] &
  \mathsf{env}_{\mathrm{defective}} \arrow[d, "\text{exaptation}"] \\
  \text{maintained in genome} &
  \mathsf{Syncytin} : \mathsf{CellFusion} \arrow[l, "\text{essential host gene}"']
\end{tikzcd}
\]
The fusogenic activity originally used for viral entry has been repurposed for
trophoblast fusion in placenta formation.
\end{example}

\begin{proposition}[Convergent Domestication]\label{prop:convergent}
Syncytin domestication has occurred independently at least 10 times across mammalian
lineages, demonstrating that:
\[
  \mathsf{naturalize} : \ERV_{\mathsf{env}} \to \mathsf{Syncytin}
\]
is a \emph{convergent morphism}---the same categorical arrow is discovered independently
by different evolutionary lineages.
\end{proposition}


% ══════════════════════════════════════════════════════════════════════
\section{The Self-Modification Paradox}\label{sec:goedel}
% ══════════════════════════════════════════════════════════════════════

We now establish a fundamental limitation on transposable elements, analogous to
G\"odel's incompleteness theorems.

\begin{theorem}[Transposon Self-Prediction Impossibility]\label{thm:goedel}
No transposable element can contain a subroutine that predicts its own next insertion site.
Formally, there is no computable function
\[
  \mathsf{predict} : \TE \times \Genome \to \mathbb{N}
\]
encoded within $\TE$ such that $\mathsf{predict}(t, g) = i$ implies $\tau(t, g)$
inserts $t$ at position $i$ in $g$.
\end{theorem}

\begin{proof}
By diagonalization. Suppose such a $\mathsf{predict}$ exists and is encoded within
the transposon $t$. Define a modified transposon $t^*$ that:
\begin{enumerate}[label=\arabic*.]
  \item Computes $i = \mathsf{predict}(t^*, g)$,
  \item Inserts itself at position $i + 1$ instead.
\end{enumerate}
Then $\mathsf{predict}(t^*, g) = i$ but $t^*$ inserts at $i + 1$, contradicting the
correctness of $\mathsf{predict}$. If we modify $\mathsf{predict}$ to account for
$t^*$, we can construct $t^{**}$, and so on. This is a direct analogue of the
G\"odelian diagonal argument: the transposon cannot simultaneously be the predictor
and the predicted.

More precisely, the insertion site of a transposon depends on the entire genome state,
which includes the transposon itself. If the transposon modifies its behavior based
on its own prediction, the genome state changes, invalidating the prediction. This
is a fixed-point obstruction: the function
\[
  F(i) = \text{``insertion site given that the transposon believes the site is } i\text{''}
\]
has no fixed point in general, by an argument analogous to Lawvere's fixed-point
theorem applied contravariantly.
\end{proof}

\begin{corollary}[Fundamental Unpredictability of Genome Evolution]\label{cor:unpredictable}
The long-term trajectory of a self-modifying genome is fundamentally unpredictable
from within, even given complete knowledge of the current genome state and all active
transposable elements.
\end{corollary}

\begin{remark}
This result does not preclude \emph{statistical} models of transposon insertion
preferences (e.g., L1 preference for AT-rich regions). It states that no
\emph{deterministic, self-referential} prediction is possible---a distinction parallel
to the difference between thermodynamic ensembles (statistical, predictable) and
individual molecular trajectories (chaotic, unpredictable).
\end{remark}

\subsection{Connection to Halting Problem}

\begin{proposition}[Transposon Halting]\label{prop:halting}
The question ``Will transposon $t$ eventually become fixed (cease transposing) in
genome $g$?'' is undecidable in general.
\end{proposition}

\begin{proof}
We reduce the halting problem to transposon fixation. Given a Turing machine $M$
with input $x$, encode $M$'s transition function as a transposase-like enzyme acting
on a genome that encodes $x$. The transposon reaches a fixed position (ceases
transposing) if and only if $M$ halts on $x$. Since halting is undecidable,
transposon fixation is undecidable.
\end{proof}


% ══════════════════════════════════════════════════════════════════════
\section{Haskell Implementation}\label{sec:implementation}
% ══════════════════════════════════════════════════════════════════════

We provide a complete Haskell implementation of the framework described above. The
implementation is organized into four modules:

\begin{itemize}
  \item \texttt{Repeat.hs}: Core type definitions for the $\Repeat\langle S \rangle$ type.
  \item \texttt{Transposon.hs}: Cut-and-paste mechanics, insertion site selection, TIR recognition.
  \item \texttt{SatelliteDNA.hs}: Tandem repeat modeling, centromeric and telomeric satellites.
  \item \texttt{DynamicGenome.hs}: Genome-level transposition simulation, composition tracking.
  \item \texttt{Main.hs}: Demonstrations and examples.
\end{itemize}

\subsection{Core Types}

The core $\Repeat$ type is implemented as a Haskell algebraic data type parameterized
by an activity state:

\begin{verbatim}
data ActivityState = Active | Suppressed | Fossilized | Domesticated

data Repeat s where
  DNATransposon   :: { tirLeft      :: Sequence
                     , tirRight     :: Sequence
                     , transposase  :: Gene
                     , dtActivity   :: s } -> Repeat s
  Retrotransposon :: { rtClass      :: RTClass
                     , reverseTranscriptase :: Maybe Gene
                     , rtActivity   :: s } -> Repeat s
  SatelliteDNA    :: { repeatUnit   :: Sequence
                     , copyNumber   :: Int
                     , satLocation  :: SatelliteLocation } -> Repeat s
  ERV             :: { ltrLeft      :: Sequence
                     , ltrRight     :: Sequence
                     , viralGenes   :: [Gene]
                     , ervIntact    :: Bool
                     , ervActivity  :: s } -> Repeat s
\end{verbatim}

\subsection{Cut-and-Paste Transposition}

The cut-and-paste operation is implemented as a pure function:

\begin{verbatim}
cutAndPaste :: Int -> Int -> Genome -> Maybe Genome
cutAndPaste from to genome
  | from < 0 || from >= length genome = Nothing
  | to   < 0 || to   >= length genome = Nothing
  | otherwise = Just (excise from >> insert to)
\end{verbatim}

\subsection{Retrotransposition}

Copy-and-paste retrotransposition creates a new copy:

\begin{verbatim}
copyAndPaste :: Int -> Int -> Int -> Genome -> Maybe Genome
copyAndPaste start len target genome =
  let element = take len (drop start genome)
  in Just (insertAt target element genome)
\end{verbatim}

\subsection{Domestication}

Transposon domestication is modeled as a type-level state change:

\begin{verbatim}
domesticate :: Repeat Active -> String -> Repeat Domesticated
domesticate (DNATransposon tl tr tp _) func =
  DNATransposon tl tr tp (Domesticated func)
\end{verbatim}

\subsection{Running the Implementation}

The implementation compiles with standard GHC:

\begin{verbatim}
$ cd src/repetitive-elements
$ ghc -o main Main.hs
$ ./main
\end{verbatim}

Output demonstrates cut-and-paste transposition, copy-and-paste retrotransposition,
satellite DNA modeling, ERV structure, and domestication events, with genome composition
statistics at each step.


% ══════════════════════════════════════════════════════════════════════
\section{Discussion and Conclusion}\label{sec:discussion}
% ══════════════════════════════════════════════════════════════════════

\subsection{Summary of Results}

We have established a comprehensive type-theoretic framework for understanding
repetitive elements as self-modifying code:

\begin{enumerate}[label=\arabic*.]
  \item \textbf{DNA transposons} implement \emph{move semantics} via cut-and-paste,
    preserving a linear type invariant (\cref{thm:linear}).
  \item \textbf{Retrotransposons} implement \emph{copy semantics}, driving monotonic
    genome growth (\cref{thm:growth}).
  \item \textbf{LINEs} are \emph{autonomous macros} with monadic structure (\cref{thm:l1-monad}).
  \item \textbf{SINEs} are \emph{dependent macros} formalized via dependent types
    (\cref{thm:sine-dep}).
  \item \textbf{Satellite DNA} provides \emph{memory alignment} at centromeres and
    telomeres (\cref{thm:centromere}).
  \item \textbf{ERVs} are \emph{legacy imports} that decay into solo LTRs with
    retained regulatory potential (\cref{thm:solo-ltr}).
  \item \textbf{Domestication} is \emph{type naturalization}, exemplified by RAG
    (\cref{thm:rag}) and syncytins (\cref{ex:syncytin}).
  \item A \textbf{G\"odelian limitation} prevents transposon self-prediction
    (\cref{thm:goedel}).
  \item \textbf{Transposition} defines an \emph{endofunctor} on the genome category
    (\cref{thm:endofunctor}).
\end{enumerate}

\subsection{Implications for Genome Engineering}

Understanding repetitive elements as self-modifying code has practical implications.
Transposon-based tools (Sleeping Beauty, piggyBac) are already used for insertional
mutagenesis and gene therapy. Our framework provides a formal basis for reasoning about:

\begin{itemize}
  \item \textbf{Insertion safety:} The endofunctor structure (\cref{thm:endofunctor})
    gives a categorical framework for analyzing how insertions compose with existing
    genomic features.
  \item \textbf{Off-target effects:} The self-prediction impossibility (\cref{thm:goedel})
    formalizes the inherent unpredictability of insertion sites.
  \item \textbf{Domestication engineering:} The naturalization framework (\cref{def:naturalization})
    suggests a principled approach to converting transposon-based tools into stable
    genomic components.
\end{itemize}

\subsection{Connections to the Modular Framework}

Within the DNA-Lang modular framework, repetitive elements interact with every other
layer:

\begin{itemize}
  \item \textbf{Paper 1 (Regulatory Sequences):} Transposon insertions create novel
    regulatory elements; solo LTRs act as enhancers and promoters.
  \item \textbf{Paper 2 (Coding Sequences):} Exon shuffling mediated by transposons
    creates new protein-coding genes.
  \item \textbf{Paper 3 (Non-Coding RNA):} SINEs generate functional non-coding RNAs;
    Alu elements are a major source of A-to-I editing substrates.
  \item \textbf{Paper 4 (Epigenetic Marks):} Epigenetic silencing (methylation,
    heterochromatin) is the primary host defense against active transposons.
  \item \textbf{Paper 6 (Developmental Programs):} Transposon-derived enhancers drive
    tissue-specific gene expression during development.
  \item \textbf{Paper 7 (Structural DNA):} Satellite DNA at centromeres and telomeres
    provides the structural scaffold of chromosomes.
\end{itemize}

\subsection{Future Directions}

Several directions remain for future work:

\begin{enumerate}[label=\arabic*.]
  \item \textbf{Stochastic transposition:} Extending the endofunctor to a \emph{probabilistic}
    endofunctor incorporating target site preferences and epigenetic accessibility.
  \item \textbf{TE ecosystem dynamics:} Modeling the competitive and cooperative
    interactions among different TE families as a categorical game.
  \item \textbf{Horizontal transfer:} Formalizing the ``import'' of new TEs from other
    species as inter-genome functors.
  \item \textbf{CRISPR as TE defense:} Connecting prokaryotic CRISPR-Cas systems to
    the TE framework as a ``type-checking firewall.''
\end{enumerate}

\subsection{Conclusion}

The view of repetitive elements as self-modifying code resolves the ``junk DNA'' paradox:
these elements are not genomic waste but the genome's metaprogramming layer. DNA transposons
implement move semantics, retrotransposons implement copy semantics, satellite DNA provides
structural alignment, ERVs are legacy imports, and domesticated elements demonstrate that
even ``parasitic'' code can become essential infrastructure. The G\"odelian self-prediction
impossibility reveals a deep connection between genome evolution and the foundations of
computation, suggesting that the creative unpredictability of evolution is not a bug but
a fundamental feature of self-modifying systems.

% ── References ────────────────────────────────────────────────────────
\begin{thebibliography}{99}

\bibitem{lander2001}
Lander, E.S., et al.
\newblock Initial sequencing and analysis of the human genome.
\newblock \emph{Nature}, 409(6822):860--921, 2001.

\bibitem{debarros2008}
de Koning, A.P., Gu, W., Castoe, T.A., Batzer, M.A., and Pollock, D.D.
\newblock Repetitive elements may comprise over two-thirds of the human genome.
\newblock \emph{PLoS Genetics}, 7(12):e1002384, 2011.

\bibitem{orgel1980}
Orgel, L.E. and Crick, F.H.
\newblock Selfish DNA: the ultimate parasite.
\newblock \emph{Nature}, 284(5757):604--607, 1980.

\bibitem{doolittle1980}
Doolittle, W.F. and Sapienza, C.
\newblock Selfish genes, the phenotype paradigm and genome evolution.
\newblock \emph{Nature}, 284(5757):601--603, 1980.

\bibitem{feschotte2008}
Feschotte, C.
\newblock Transposable elements and the evolution of regulatory networks.
\newblock \emph{Nature Reviews Genetics}, 9(5):397--405, 2008.

\bibitem{kazazian2004}
Kazazian, H.H.
\newblock Mobile elements: drivers of genome evolution.
\newblock \emph{Science}, 303(5664):1626--1632, 2004.

\bibitem{agrawal1998}
Agrawal, A., Eastman, Q.M., and Schatz, D.G.
\newblock Transposition mediated by RAG1 and RAG2 and its implications for the
evolution of the immune system.
\newblock \emph{Nature}, 394(6695):744--751, 1998.

\bibitem{hiom1998}
Hiom, K., Melek, M., and Gellert, M.
\newblock DNA transposition by the RAG1 and RAG2 proteins: a possible source
of oncogenic translocations.
\newblock \emph{Cell}, 94(4):463--470, 1998.

\bibitem{sanmiguel1998}
SanMiguel, P., Gaut, B.S., Tikhonov, A., Nakajima, Y., and Bennetzen, J.L.
\newblock The paleontology of intergene retrotransposons of maize.
\newblock \emph{Nature Genetics}, 20(1):43--45, 1998.

\bibitem{mi2000}
Mi, S., Lee, X., Li, X., Veldman, G.M., Bhatt, R.S., et al.
\newblock Syncytin is a captive retroviral envelope protein involved in human
placental morphogenesis.
\newblock \emph{Nature}, 403(6771):785--789, 2000.

\bibitem{wadsworth2009}
Wadsworth, C. and Bhatt, P.
\newblock Modular composition in biological systems: from molecular to organismal scales.
\newblock \emph{Journal of Theoretical Biology}, 261(3):401--412, 2009.

\bibitem{awodey2010}
Awodey, S.
\newblock \emph{Category Theory}.
\newblock Oxford University Press, 2nd edition, 2010.

\bibitem{pierce2002}
Pierce, B.C.
\newblock \emph{Types and Programming Languages}.
\newblock MIT Press, 2002.

\bibitem{maclane1998}
Mac~Lane, S.
\newblock \emph{Categories for the Working Mathematician}.
\newblock Springer, 2nd edition, 1998.

\bibitem{lawvere1969}
Lawvere, F.W.
\newblock Diagonal arguments and cartesian closed categories.
\newblock \emph{Lecture Notes in Mathematics}, 92:134--145, 1969.

\bibitem{thompson2016}
Thompson, P.J., Macfarlan, T.S., and Lorincz, M.C.
\newblock Long terminal repeats: from parasitic elements to building blocks of
the transcriptional regulatory repertoire.
\newblock \emph{Molecular Cell}, 62(5):766--776, 2016.

\bibitem{bourque2018}
Bourque, G., Burns, K.H., Gehring, M., et al.
\newblock Ten things you should know about transposable elements.
\newblock \emph{Genome Biology}, 19:199, 2018.

\end{thebibliography}

\end{document}
